\documentclass[a4paper,11pt]{article}
\usepackage{exam-style-341}
\examinehead{341 农业综合知识三 · 网络技术分章练习}{第04章 参考答案与解析}
\begin{document}

\unitanswerscover{341 农业综合知识三 · 网络技术分章练习}{第04章参考答案}{数据链路层：逐题答案与解析；CRC 与滑动窗口计算题给出完整步骤，停止等待／连续 ARQ 的信道利用率给出公式。}

\answerpart{一、单项选择题（10 题，每题 3 分，共 30 分）}

\q B。\begin{analysis}数据链路层把网络层交下来的分组封装成帧，以帧为单位传输，其 PDU 是帧（教材 3.1.1，印刷页 72）。比特属于物理层，分组属于网络层，报文属于应用层。\end{analysis}

\q A。\begin{analysis}数据链路层的三个基本问题是封装成帧、透明传输和差错检测（教材 3.1.2，印刷页 73—77）。B 属于网络层和运输层功能，C 属于物理层，D 是可靠传输协议的内容。\end{analysis}

\q A。\begin{analysis}PPP 帧以标志字段 0x7E（01111110）作为帧的定界符（教材 3.2.2，印刷页 80—82）。0x7D 是转义符，0xFF 和 0x03 是地址字段和控制字段的固定值。\end{analysis}

\q B。\begin{analysis}PPP 使用异步传输时采用字节填充，把转义符定义为 0x7D：数据中的 0x7E 变为 0x7D 0x5E，0x7D 变为 0x7D 0x5D（教材 3.1.2，印刷页 73—75）。\end{analysis}

\q A。\begin{analysis}零比特填充的规则是：发送端扫描信息字段，只要发现 5 个连续 1，就立即填入 1 个 0，从而保证信息字段中不会出现 6 个连续 1（与标志 01111110 区分）；接收端每发现 5 个连续 1 就把其后的 0 删除（教材 3.1.2，印刷页 75）。\end{analysis}

\q A。\begin{analysis}CRC 采用模 2 除法：把数据乘以 $x^n$（末尾补 $n$ 个 0）后除以生成多项式，所得余数即为帧检验序列 FCS（教材 3.1.2，印刷页 75—77）。\end{analysis}

\q A。\begin{analysis}FCS 用于差错检测：接收端用同样的生成多项式去除收到的帧，若余数为 0 则认为传输无差错（并非绝对可靠，但漏检概率极低）。FCS 本身不能纠正差错，纠错要靠重传（教材 3.1.2）。\end{analysis}

\q A。\begin{analysis}停止等待协议中发送方每发一帧就等待确认，接收方只需区分“新帧”和“重传帧”，因此 1 位序号（0 和 1 交替）即可；这与提高速率无关（教材第 8 版未展开，按通用教材补充）。\end{analysis}

\q A。\begin{analysis}回退 N 帧（连续 ARQ）的接收窗口为 1：接收方按序接收，某帧出错后丢弃该帧及其后所有帧，发送方从出错帧开始重传（教材第 8 版未展开，按通用教材补充）。\end{analysis}

\q A。\begin{analysis}PPP 协议由三部分组成：一个将 IP 数据报封装到串行链路的方法（封装成帧）、链路控制协议 LCP、网络控制协议 NCP（教材 3.2.1，印刷页 78—79）。\end{analysis}

\answerpart{二、多项选择题（3 题，每题 4 分，共 12 分）}

\q ABC。\begin{analysis}数据链路层的三个基本问题是封装成帧、透明传输和差错检测（教材 3.1.2）。D 路由选择属于网络层功能。\end{analysis}

\q ABC。\begin{analysis}PPP 的组成部分为封装成帧、LCP 和 NCP（教材 3.2.1）。D 路由选择协议 RIP 是网络层协议，不属于 PPP 的组成部分。\end{analysis}

\q ABC。\begin{analysis}A 正确：收发双方必须使用相同的生成多项式；B 正确：CRC 计算采用模 2 除法；C 正确：CRC 可以检测出多位比特差错；D 错误：CRC 只能检错，不能自动纠正差错，纠错需依靠重传（教材 3.1.2）。\end{analysis}

\answerpart{三、判断题（6 题，每题 2 分，共 12 分）}

\q ×。改正：链路是一段无源的点到点物理线路段，数据链路则是链路加上实现通信协议的硬件和软件。\begin{analysis}教材 3.1.1（印刷页 72）区分了链路与数据链路：链路只提供物理传输通路，数据链路还要有协议软硬件来保证帧的正确传输。\end{analysis}

\q √。\begin{analysis}透明传输是指无论数据部分出现什么比特组合，都能按原样正确传输，不会与帧定界符混淆；PPP 用字节填充或零比特填充实现透明传输（教材 3.1.2，印刷页 73—75）。\end{analysis}

\q ×。改正：CRC 循环冗余检验只能检测出差错，不能自动纠正差错，出错后通常依靠重传恢复。\begin{analysis}CRC 的漏检概率很低，但接收端只能判断“收到的帧是否有错”，不能定位并纠正错误比特（教材 3.1.2，印刷页 75—77）。\end{analysis}

\q √。\begin{analysis}停止等待协议中发送方每发送一个帧就停止发送并等待确认，超时未收到确认就重传该帧（教材第 8 版未展开，按通用教材补充）。\end{analysis}

\q ×。改正：回退 N 帧协议的接收窗口为 1，选择重传 ARQ 的接收窗口才大于 1。\begin{analysis}回退 N 帧按序接收，接收方只期待当前序号；某帧出错后，该帧及其后的帧都要重传（通用教材）。\end{analysis}

\q √。\begin{analysis}PPP 在异步传输时使用字节填充，在同步传输（如 SONET/SDH）时使用零比特填充，两者都能实现透明传输（教材 3.1.2、3.2.2）。\end{analysis}

\answerpart{四、计算题（2 题，每题 5 分，共 10 分）}

\q 参考答案：
（1）生成多项式 $G(x)=x^3+x+1$，除数 1011，最高次为 3，故在数据 101110 后补 3 个 0，得被除数 101110000（1 分）。
用模 2 除法（多项式除法）：
数据多项式 $M(x)=x^5+x^3+x^2+x$，乘 $x^3$ 得 $x^8+x^6+x^5+x^4$；
除以 $G(x)$：$x^8+x^6+x^5+x^4=(x^5)(x^3+x+1)+x^4$，余 $x^4$；
$x^4=(x)(x^3+x+1)+(x^2+x)$，余 $x^2+x$，即二进制 110（3 分）。
所以 FCS $=110$。
（2）实际发送的比特序列为数据后接 FCS：$101110\,110$（1 分）。
\begin{analysis}校验方法：接收端用 1011 去除收到的 101110110，若余数为 0 则认为无差错。注意 CRC 的“除法”是模 2 除法，加减法均为异或，不产生借位和进位。教材 3.1.2 的例 3-1 与本題方法相同。\end{analysis}

\q 参考答案：
（1）停止等待协议：发送窗口 $=1$，接收窗口 $=1$（1 分，1 位序号即可，本题序号为 3 位是冗余的）。
（2）连续 ARQ（回退 N 帧）：发送窗口 $\leq 2^n-1=2^3-1=7$，接收窗口 $=1$（2 分）。
（3）选择重传 ARQ：发送窗口 $\leq 2^{n-1}=2^{3-1}=4$，接收窗口 $=4$（2 分）。
\begin{analysis}窗口大小的限制目的都是防止接收方把新帧误认为重传帧（或相反）。回退 N 帧要求“发送窗口 $+$ 接收窗口 $\leq 2^n$”，因接收窗口为 1，故发送窗口最大 $2^n-1$；选择重传要求收发窗口均不超过 $2^{n-1}$，否则序号会产生歧义。（教材第 8 版未展开，按通用教材补充）\end{analysis}

\answerpart{五、简答题（2 题，每题 8 分，共 16 分）}

\q 参考答案要点：
（1）三个基本问题（3 分）：封装成帧——在数据前后加首部和尾部，使接收端能识别帧的边界；透明传输——无论数据出现什么比特组合都能正确传输，不会与帧定界符混淆；差错检测——用 CRC 等校验方法发现传输中的比特差错（教材 3.1.2，印刷页 73—77）。
（2）透明传输的两种实现（5 分）：（1） 字节填充：适用于异步传输，PPP 把转义符定义为 0x7D，数据中的 0x7E 变为 0x7D 0x5E，0x7D 变为 0x7D 0x5D，控制字符前也加 0x7D 并改变编码；（2） 零比特填充：适用于同步传输，发送端每发现 5 个连续 1 就填入 1 个 0，接收端每发现 5 个连续 1 就删除其后的 0。二者分别适用于“逐字符传送”和“连续比特流传送”的场合。

\q 参考答案要点：
（1）三个组成部分（3 分）：PPP 由封装成帧、链路控制协议 LCP 和网络控制协议 NCP 三部分组成；LCP 用于建立、配置、测试和释放数据链路连接，NCP 用于配置网络层协议（如为接入主机分配 IP 地址）（教材 3.2.1，印刷页 78—79）。
（2）帧格式主要字段（3 分）：PPP 帧以标志字段 0x7E 定界，依次包含地址字段 0xFF、控制字段 0x03、协议字段（2 B，指明信息字段中的协议类型，如 0x0021 表示 IP）、信息字段（长度可变）、帧检验序列 FCS（2 B），最后再以标志字段 0x7E 结束（教材 3.2.2，印刷页 80—82）。
（3）透明传输的实现（2 分）：异步传输时用字节填充（转义符 0x7D）；同步传输时用零比特填充（5 个连续 1 后填 0）。这样即使信息字段中出现与标志字段相同的比特组合，也不会被误认为帧边界。

\answerpart{六、综合题（1 题，共 20 分）}

\q 参考答案要点：
（1）（4 分）数据帧发送时延 $T_D=1000\ \text{bit}\div(10\times10^{6}\ \text{bit/s})=1\times10^{-4}\ \text{s}=0.1\ \text{ms}=100\ \mu\text{s}$；单向传播时延 $T_{\text{prop}}=10\ \text{ms}$；确认帧发送时延 $T_A=100\ \text{bit}\div(10\times10^{6}\ \text{bit/s})=1\times10^{-5}\ \text{s}=0.01\ \text{ms}=10\ \mu\text{s}$。
（2）（6 分）停止等待协议发送一帧并收到确认的周期：
$T=T_D+T_{\text{prop}}+T_A+T_{\text{prop}}=0.1+10+0.01+10=20.11\ \text{ms}$；
信道利用率 $U=\dfrac{T_D}{T}=\dfrac{0.1}{20.11}\approx0.00497\approx0.5\%$。
（3）（5 分）连续 ARQ（回退 N 帧），发送窗口 $W=4$，忽略确认帧发送时延：
$U=\dfrac{W\cdot T_D}{T_D+2T_{\text{prop}}}=\dfrac{4\times0.1}{0.1+20}=\dfrac{0.4}{20.1}\approx0.0199\approx2.0\%$。
（4）（5 分）提高信道利用率的途径：增大发送窗口、增大帧长（减少每帧的相对开销）、缩短传播时延（缩短链路距离）、提高链路数据率、采用连续 ARQ／滑动窗口代替停止等待。长时延链路上停止等待协议效率很低，因为每发送一帧后发送方必须空等一个往返时间 $2T_{\text{prop}}=20\ \text{ms}$，而真正用于发送数据的时间只有 $0.1\ \text{ms}$，等待时间远大于发送时间；连续 ARQ 可以在等待确认期间连续发送多个帧，使链路“管道”被填满，从而提高利用率。
\begin{analysis}本题是数据链路层可靠传输与信道利用率的经典计算。注意三点：一是停止等待的周期包括数据帧发送、数据传播、确认帧发送、确认传播四段；二是回退 N 帧的利用率公式仅当 $W\cdot T_D<T_D+2T_{\text{prop}}$（窗口未填满管道）时成立，本题窗口小，公式适用；三是“忽略确认帧”时把 $T_A$ 置零。\end{analysis}

\end{document}
